% SGA 3, Expose VI_B, Loop 1 English reconstruction.
% Scope: Examples 6.2.4(a)--(e) and Corollary 6.2.5 in full,
% including source editor's notes 56--58.
% Authority: Polo--Gille Exp6B-13oct24.pdf, component-local pages
% 41--43 (combined-reader pages 382--384; displayed re-edition
% headers 371--373). Corollary 6.2.6 begins on local page 43 and
% is excluded.
% OCR and Reinhold English Markdown are comparison witnesses only.
% The displayed constructions are reconstructed natively in TeX;
% this unit contains no raster diagrams.

\medskip
\phantomsection\label{ex:sga3-VIB-6.2.4}%
\noindent\textbf{例 6.2.4.---}
下面给出一些重要的函子例子, 它们可化归为
\hyperref[thm:sga3-VIB-6.2.3]{6.2.3} 中所考虑的类型的函子
\[
  \prod_{Z/S}Y/Z
\]
且后文将用到它们的可表性判据.
用 \(S\) 表示一个概形,
\(X,Y,Z,\ldots\) 表示 \(S\) 上的概形.

\medskip
\phantomsection\label{ex:sga3-VIB-6.2.4-a}%
\noindent\textbf{a)}
假设给定一个 \(S\)-态射
\begin{equation}
  q:X\longrightarrow\underline{\operatorname{Hom}}_S(Y,Z),
  \tag{x}\label{eq:sga3-VIB-6.2.4-x}
\end{equation}
(``\(X\) 作用于 \(Y\), 取值于 \(Z\)''), 即一个
态射
\begin{equation}
  r:X\times_S Y\longrightarrow Z.
  \tag{xx}\label{eq:sga3-VIB-6.2.4-xx}
\end{equation}
考虑 \(Z\) 的一个子概形 \(Z'\), 从而有单态射
\[
  \underline{\operatorname{Hom}}_S(Y,Z')
  \longrightarrow
  \underline{\operatorname{Hom}}_S(Y,Z),
\]
它将第一个函子视为第二个的子函子. 设
\(X'\) 为此子函子在 \(q\) 下的逆像. 因此
\(X'\) 为 \(X\) 的子函子, 其中 \(X'(T)\) 是满足
\[
  q(x):Y_T\longrightarrow Z_T
\]
经由 \(Z'_T\) 分解的 \(x\in X(T)\) 所成的集合.

此函子 \(X'\) 可描述如下. 令
\[
  P=X\times_S Y,
\]
并设 \(P'\) 为 \(Z'\) 在
\(r:P\to Z\) 下的逆像. 则有显然的同构
\begin{equation}
  X'\simeq\prod_{P/X}P'/P.
  \tag{xxx}\label{eq:sga3-VIB-6.2.4-xxx}
\end{equation}
因此, 若 \(Y\) 在 \(S\) 上本质自由,
且 \(Z'\) 在 \(Z\) 中闭, 则 \(X\) 的子函子 \(X'\)
可由 \(X\) 的闭子概形表出.%
\footnote{\label{note:sga3-VIB-56}编者注: 我们补充了
下面一句.}
若此外 \(Z'\to Z\) 有限呈示, 则
\(X'\to X\) 也有限呈示.

\medskip
\phantomsection\label{ex:sga3-VIB-6.2.4-b}%
\noindent\textbf{b)}
假设给定 \(X\) 作用于 \(Y\) 且取值于
\(Z\) 的两种方式, 即两个态射
\[
  q_1,q_2:X\rightrightarrows
  \underline{\operatorname{Hom}}_S(Y,Z),
\]
并令
\[
  X'=\operatorname{Ker}(q_1,q_2).
\]
因此 \(X'\) 为 \(X\) 的子函子, 其中 \(X'(T)\) 是满足两个态射
\[
  q_1(x),q_2(x):Y_T\rightrightarrows Z_T
\]
相等的 \(x\in X(T)\) 所成的集合. 给出 \(q_1,q_2\) 等价于给出
一个态射
\[
  q:X\longrightarrow
  \underline{\operatorname{Hom}}_S(Y,Z\times_S Z),
\]
或者等价地, 给出一个态射
\[
  r:X\times_S Y\longrightarrow Z\times_S Z.
\]
令 \(U=Z\times_S Z\), 并设 \(U'\) 为 \(Z\times_S Z\)
的对角子概形. 则 \(X'\) 恰为子函子
\[
  \underline{\operatorname{Hom}}_S(Y,U')
  \subset
  \underline{\operatorname{Hom}}_S(Y,U).
\]
在 \(q\) 下的逆像. 因此可以将它写成
\hyperref[eq:sga3-VIB-6.2.4-xxx]{\textup{(xxx)}} 的形式, 其中
\(P=X\times_S Y\), \(P'\) 为对角线
在 \(r\) 下的逆像, 即
\[
  r_1,r_2:X\times_S Y\rightrightarrows Z.
\]
的核. 于是处于
\hyperref[ex:sga3-VIB-6.2.4-a]{\textup{(a)}} 的假设之下.

因此, 若 \(Y\) 在 \(S\) 上本质自由, 且 \(Z\)
在 \(S\) 上分离, 则 \(X\) 的子函子 \(X'\)
可由 \(X\) 的闭子概形表出.%
\footnotemark[56]
若此外 \(Z\to S\) 局部有限型, 则
\(X'\to X\) 有限呈示.

\medskip
\phantomsection\label{ex:sga3-VIB-6.2.4-c}%
\noindent\textbf{c)}
假设给定一个态射
\[
  q:X\longrightarrow\underline{\operatorname{Hom}}_S(Y,Y),
\]
即 ``\(X\) 作用于 \(Y\).'' 设 \(X'\) 为
此态射的 ``核'': \(X'\) 是 \(X\) 的子函子,
其中 \(X'(T)\) 为满足
\[
  q(x):Y_T\longrightarrow Y_T
\]
为恒等态射的 \(x\in X(T)\) 所成的集合. 令 \(X\) 平凡作用于 \(Y\),
从而引入第二个态射
\[
  q':X\longrightarrow\underline{\operatorname{Hom}}_S(Y,Y)
\]
即可看出此函子属于 \hyperref[ex:sga3-VIB-6.2.4-b]{\textup{(b)}}
所论情形. 因此, 若 \(Y\) 在 \(S\) 上本质自由且分离,
则子函子 \(X'\), 即 \(q\) 的核, 可由
\(X\) 的闭子概形表出.%
\footnotemark[56]
若此外 \(Y\to S\) 局部有限型, 则
\(X'\to X\) 有限呈示.

\medskip
\phantomsection\label{ex:sga3-VIB-6.2.4-d}%
\noindent\textbf{d)}
在 \hyperref[ex:sga3-VIB-6.2.4-c]{(c)} 的假设下,
考虑 \(Y\) 的 ``\(X\) 下不变量'' 子函子 \(Y'\).
因此 \(Y'(T)\) 是满足如下条件的 \(y\in Y(T)\) 所成的集合:
相应的态射
\[
  \bar q(y):X_T\longrightarrow Y_T
\]
是取值为 \(y\) 的常值 \(T\)-态射.
如 \hyperref[ex:sga3-VIB-6.2.4-c]{(c)} 中一样引入
\(\bar q'\), 以及对应于 \(\bar q\) 和 \(\bar q'\) 的态射
\[
  \bar q,\bar q':Y\rightrightarrows
  \underline{\operatorname{Hom}}_S(X,Y),
\]
可知
\[
  Y'=\operatorname{Ker}(\bar q,\bar q').
\]
因此它仍属于 \hyperref[ex:sga3-VIB-6.2.4-b]{\textup{(b)}}
所论情形, 只需交换
\(X\) 和 \(Y\) 的角色, 并取 \(Z=Y\).

因此, 若 \(X\) 在 \(S\) 上本质自由, 且 \(Y\)
在 \(S\) 上分离, 则 \(Y\) 的 \(X\) 下不变量
子函子 \(Y'\) 可由 \(Y\) 的闭子概形表出.%
\footnotemark[56]
若此外 \(Y\to S\) 局部有限型, 则
\(Y'\to Y\) 有限呈示.

\medskip
\phantomsection\label{ex:sga3-VIB-6.2.4-e}%
\noindent\textbf{e)}
上述各例中描述的这类构造在群论中
尤其常见. 例如, 当 \(G\) 为一个作用于 \(S\)-概形 \(X\)
的 \(S\)-群概形时,
\[
  q:G\longrightarrow\underline{\operatorname{Aut}}_S(X),
\]
只要 \(X\) 在 \(S\) 上本质自由且分离,
\(q\) 的核, 即平凡作用的 \(G\) 的子群, 就是
\(G\) 的一个闭子概形
(\hyperref[ex:sga3-VIB-6.2.4-c]{例~(c)}). 只要
\(G\) 在 \(S\) 上本质自由且 \(X\) 在 \(S\) 上分离,
不变量子对象 \(X^G\) 就是 \(X\) 的一个闭子概形
(\hyperref[ex:sga3-VIB-6.2.4-d]{例~(d)}).

设 \(Y,Z\) 为 \(X\) 的子概形, 并考虑 \(G\) 的子函子
\(\underline{\operatorname{Transp}}_G(Y,Z)\), 即将 \(Y\)
传送到 \(Z\) 的传送子. 对于 \(S\) 上的 \(T\),
它的 \(T\)-值点是满足相应的 \(X_T\) 的自同构使
\[
  g(Y_T)\subset Z_T,
\]
成立的元素 \(g\in G(T)\), 即该自同构诱导一个
经由 \(Y_T\to Z_T\) 分解的态射 \(Y_T\to X_T\).
因此, 若 \(Y\) 在 \(S\) 上本质自由,
且 \(Z\) 在 \(X\) 中闭, 则
\(\underline{\operatorname{Transp}}_G(Y,Z)\) 为 \(G\) 的闭
子概形
(\hyperref[ex:sga3-VIB-6.2.4-a]{例~(a)}).

还可以考虑将 \(Y\) 传送到
\(Z\) 的严格传送子,%
\footnote{\label{note:sga3-VIB-57}编者注: 记为
\(\underline{\operatorname{Transp}}^{\mathrm{str}}_G(Y,Z)\).}
其在 \(S\) 上的 \(T\) 的 \(T\)-值点是满足
\(g(Y_T)=Z_T\) 的元素 \(g\in G(T)\). 此函子恰为
\[
  \underline{\operatorname{Transp}}_G(Y,Z)
  \cap
  \sigma\bigl(\underline{\operatorname{Transp}}_G(Z,Y)\bigr),
\]
其中 \(\sigma\) 为 \(G\) 的取逆态射. 因此,
若 \(Y\) 和 \(Z\) 在 \(S\) 上本质自由且在 \(X\)
中闭, 则将 \(Y\) 传送到 \(Z\) 的严格传送子为 \(G\) 的闭
子概形.

一个重要情形是 \(X=G\), 且 \(G\) 通过内自同构
作用于自身. 若 \(H\) 为 \(G\) 的子概形, 则将 \(H\) 传送到 \(H\) 的严格
传送子也称为 \emph{\(H\) 在 \(G\) 中的正规化子}, 记为
\(\underline{\operatorname{Norm}}_G H\). 因此, 若 \(H\) 为
\(G\) 的闭子群概形, 且在 \(S\) 上本质自由,
则 \(\underline{\operatorname{Norm}}_G H\) 可由
\(G\) 的闭子群概形表出.

最后, 设 \(Z\) 为 \(G\) 的子概形. 它在 \(G\) 中的中心化子
\(\underline{\operatorname{Centr}}_G(Z)\) 为 \(G\) 的群
子函子, 当将 \(Z\) 视为通过 \(G\) 的运算所诱导的运算
作用于 \(G\) 时, 它由
\hyperref[ex:sga3-VIB-6.2.4-d]{(d)} 的过程定义.
因此, 若 \(Z\) 在 \(S\) 上本质自由, 且 \(G\) 在
\(S\) 上分离, 则
\(\underline{\operatorname{Centr}}_G(Z)\) 为 \(G\) 的闭子群
概形. 特别地, 若 \(G\) 在 \(S\) 上本质自由
且分离, 则 \(G\) 的中心 \(C\), 即
\(\underline{\operatorname{Centr}}_G(G)\), 为
\(G\) 的闭子群概形.

当 \(S\) 为一个域的谱时,
\hyperref[prop:sga3-VIB-6.2.2]{6.2.2\textup{(b)}} 表明,
上述各例
\hyperref[ex:sga3-VIB-6.2.4-a]{(a)}--%
\hyperref[ex:sga3-VIB-6.2.4-e]{(e)} 中的
``本质自由'' 条件均自动满足, 只剩下分离性条件.
回顾域上的群概形必然分离
(\hyperref[sga3:VIA:0.3]{VI\(_A\), 0.3}), 例如可得:

\medskip
\phantomsection\label{cor:sga3-VIB-6.2.5}%
\noindent\textbf{推论 6.2.5.---}
\emph{设 \(G\) 为域 \(k\) 上的群概形,
\(Y,Y'\) 为 \(G\) 的两个子概形. 则:}
\begin{enumerate}
\item[\textup{(i)}]\phantomsection\label{concl:sga3-VIB-6.2.5-i}
  \emph{\(Y\) 在 \(G\) 中的中心化子为 \(G\) 的闭子群
  概形; 特别地, \(G\) 的中心
  \(\underline{\operatorname{Centr}}_G(G)\) 具有这一性质.}
\item[\textup{(i\(^{\prime}\))}]
  \phantomsection\label{concl:sga3-VIB-6.2.5-iprime}
  \emph{更一般地, 若 \(u,v:X\to G\) 为概形
  态射, 则 \(\underline{\operatorname{Transp}}_G(u,v)\)
  可由 \(G\) 的闭子概形表出.}
\item[\textup{(ii)}]\phantomsection\label{concl:sga3-VIB-6.2.5-ii}
  \emph{若 \(Y\) 闭, 则传送子
  \(\underline{\operatorname{Transp}}_G(Y',Y)\) 为 \(G\) 的闭
  子概形. 若 \(Y'\) 也闭, 则同一结论对
  \(\underline{\operatorname{Transp}}^{\mathrm{str}}_G(Y',Y)\)
  也成立.}
\item[\textup{(iii)}]\phantomsection\label{concl:sga3-VIB-6.2.5-iii}
  \emph{对 \(G\) 的每个子群概形}%
  \footnote{\label{note:sga3-VIB-58}编者注: 事实上, 在
  域 \(k\) 上, \(G\) 的每个子群概形均闭; 见
  \hyperref[sga3:VIA:0.5.2]{VI\(_A\), 0.5.2}. 此外, \SGARefLink{cor:sga3-VIB-6.2.5}{6.2.5}
  完成了从 VIII, \S6 摘取的结果的插入.}
  \emph{\(H\),
  \(\underline{\operatorname{Norm}}_G(H)\) 为 \(G\) 的闭子群
  概形.}
\end{enumerate}

% 翻译统计: 源行数 294; 输出行数 289 (含本统计行); 行内/行间数学片段 225; 翻译占位 0. 底本疑点另报主线程.

% 独立复核检查点：已逐段对照英文源文，覆盖正文、假设、证明、脚注与源注。
% 行内/陈列数学及全部引用键机械对照通过；仅统一群的invariant为“正规”，保留源文疑点另报。
